% !TeX program = lualatex
% !TeX encoding = utf8
% !TeX spellcheck = uk_UA
% !TeX root = main.tex

%% --------------------------------------------------------
%% Додаток: таблиця позначень (поки що Basics, Waves і розд. «Хвилі в замагніченій плазмі»)
%% --------------------------------------------------------

\section{Таблиця позначень}\label{sec:notation}

Нижче зібрано позначення, що вживаються в~розд.~\ref{sec:plasma-basics}, \ref{sec:plasma-waves} і~\ref{sec:mhd-waves}, у~порядку
їхнього введення. Усі формули подано в~системі СГС (гаусовій), як у~розд.~\ref{sec:plasma-basics}; температури $T$
(у~тому числі $T_e$, $T_i$) записано в~енергетичних одиницях, $k_{\mathrm B}=1$, а~числові значення дають
в~електронвольтах: $1~\text{еВ}\approx1{,}16\times10^{4}$~К. Одиниці СГС потенціалу й~поля відповідають
$300$~В і~$300$~В$/$см. Стовпець «Де» вказує
параграф або формулу, де величину введено; символи, що мають кілька значень залежно від контексту,
перелічено в~кінці додатка.

{\small
\DefTblrTemplate{caption}{default}{}% без підпису «Таблиця А.N»
\DefTblrTemplate{conthead}{default}{}
\DefTblrTemplate{capcont}{default}{}
\DefTblrTemplate{firsthead}{default}{}
\DefTblrTemplate{middlehead}{default}{}
\DefTblrTemplate{lasthead}{default}{}
\DefTblrTemplate{contfoot}{default}{}% без «Продовження на наступній сторінці»
\par\Needspace{7\baselineskip}\medskip\noindent\textbf{Загальні сталі та позначення}\par\nobreak\smallskip
\noindent\begin{longtblr}{colspec={X[0.9,c,m]X[3.4,l,m]X[1.5,c,m]X[1.5,c,m]}}
Символ & Значення & Одиниця (СГС) & Де \\
$e$ & елементарний заряд & $4{,}80\times10^{-10}$~СГСЕ & §\,\ref{sec:debye-screening}\,(с.\,\pageref{sec:debye-screening}) \\
$m_e$, $m_i$, $m_a$, $m$ & маса електрона, іона, атома; $m$~--- маса частинки (у~формулі Саха~--- електрона) & г & §\,\ref{sec:saha}\,(с.\,\pageref{sec:saha}) \\
$k_{\mathrm B}$ & стала Больцмана (далі $k_{\mathrm B}=1$) & ерг/К & §\,\ref{sec:what-is-plasma}\,(с.\,\pageref{sec:what-is-plasma}) \\
$\hbar$, $h=2\pi\hbar$ & стала Планка & ерг$\cdot$с & \eqref{eq:saha} \\
$c$ & швидкість світла & см/с & \eqref{eq:larmor} \\
$Z$ & кратність іонізації: заряд іона $Ze$ & --- & \eqref{eq:quasineutrality} \\
$n$, $n_e$, $n_i$, $n_n$ & концентрація (загальна), електронів, іонів, нейтральних частинок & см$^{-3}$ & \eqref{eq:ionization-degree} \\
$T$, $T_e$, $T_i$, $T_s$ & температура (в~енергетичних одиницях): загальна, електронів, іонів, сорту $s$ & еВ & §\,\ref{sec:two-temperatures}\,(с.\,\pageref{sec:two-temperatures}) \\
\end{longtblr}

\par\Needspace{7\baselineskip}\medskip\noindent\textbf{Іонізація (формула Саха)}\par\nobreak\smallskip
\noindent\begin{longtblr}{colspec={X[0.9,c,m]X[3.4,l,m]X[1.5,c,m]X[1.5,c,m]}}
Символ & Значення & Одиниця (СГС) & Де \\
$\alpha$ & ступінь іонізації $n_i/(n_i+n_n)$ & --- & \eqref{eq:ionization-degree} \\
$I$ & енергія іонізації атома (для водню $13{,}6$~еВ) & еВ & \eqref{eq:saha} \\
$g$, $g_e$, $g_i$, $g_a$ & статистичні ваги & --- & §\,\ref{sec:saha}\,(с.\,\pageref{sec:saha}) \\
$\mu$ & хімічний потенціал & еВ & §\,\ref{sec:saha}\,(с.\,\pageref{sec:saha}) \\
$\lambda_T$, $\lambda_{Te}$ & теплова довжина хвилі де~Бройля $h/\sqrt{2\pi mT}$ (для електрона~--- $\lambda_{Te}$) & см & §\,\ref{sec:saha}\,(с.\,\pageref{sec:saha}) \\
$n_q$ & квантова концентрація $\lambda_{Te}^{-3}=(mT/2\pi\hbar^2)^{3/2}$ & см$^{-3}$ & §\,\ref{sec:saha}\,(с.\,\pageref{sec:saha}) \\
$\Delta E/E$ & частка енергії, яку електрон передає іону при пружному зіткненні & --- & §\,\ref{sec:two-temperatures}\,(с.\,\pageref{sec:two-temperatures}) \\
\end{longtblr}

\par\Needspace{7\baselineskip}\medskip\noindent\textbf{Квазінейтральність і дебаївське екранування}\par\nobreak\smallskip
\noindent\begin{longtblr}{colspec={X[0.9,c,m]X[3.4,l,m]X[1.5,c,m]X[1.5,c,m]}}
Символ & Значення & Одиниця (СГС) & Де \\
$\rho$ & густина заряду $\sum_sq_sn_s$ & СГСЕ$/$см$^{3}$ & \eqref{eq:quasineutrality} \\
$\vect E$, $E$ & напруженість електричного поля & СГС (поле) & §\,\ref{sec:quasineutrality}\,(с.\,\pageref{sec:quasineutrality}) \\
$\delta n$ & локальне відхилення концентрації від середньої & см$^{-3}$ & §\,\ref{sec:quasineutrality}\,(с.\,\pageref{sec:quasineutrality}) \\
$L$ & розмір плазми (характерний масштаб неоднорідності) & см & \eqref{eq:size-condition} \\
$\varphi$ & електростатичний потенціал; $\varphi_0=\varphi(0)$~--- потенціал поверхні & СГС (потенціал) & §\,\ref{sec:debye-screening}\,(с.\,\pageref{sec:debye-screening}) \\
$q$, $q_s$ & пробний заряд; заряд частинки сорту $s$ & СГСЕ & §\,\ref{sec:debye-screening}\,(с.\,\pageref{sec:debye-screening}) \\
$Q$ & густина потоку енергії на стінку (на відміну від заряду $q$) & ерг$\cdot$см$^{-2}\cdot$с$^{-1}$ & §\,\ref{sec:bohm-sheath}\,(с.\,\pageref{sec:bohm-sheath}) \\
$L_\mathrm{ps}$ & масштаб передшару (на відміну від розміру плазми $L$) & см & §\,\ref{sec:bohm-sheath}\,(с.\,\pageref{sec:bohm-sheath}) \\
$n_{s0}$ & незбурена концентрація сорту $s$ & см$^{-3}$ & §\,\ref{sec:debye-screening}\,(с.\,\pageref{sec:debye-screening}) \\
$\lambda_D$, $\lambda_{De}$, $\lambda_{Di}$ & дебаївська довжина (повна, електронна, іонна); $\lambda_D=\lambda_{De}/\sqrt2$ при $T_e=T_i$ & см & \eqref{eq:poisson-screened}, \eqref{eq:debye-radius} \\
$\delta(\vect r)$ & дельта-функція (точковий заряд) & см$^{-3}$ & \eqref{eq:poisson-screened} \\
\end{longtblr}

\par\Needspace{7\baselineskip}\medskip\noindent\textbf{Екранування у виродженому газі}\par\nobreak\smallskip
\noindent\begin{longtblr}{colspec={X[0.9,c,m]X[3.4,l,m]X[1.5,c,m]X[1.5,c,m]}}
Символ & Значення & Одиниця (СГС) & Де \\
$E_F$ & енергія Фермі & еВ & \eqref{eq:thomas-fermi} \\
$\lambda_{TF}$ & довжина Томаса--Фермі & см & \eqref{eq:thomas-fermi} \\
$\lambda$ & довжина екранування в~загальному випадку $[4\pi e^2\partial n/\partial\mu]^{-1/2}$ & см & \eqref{eq:screening-general} \\
$a_B$ & борівський радіус $\hbar^2/me^2=0{,}53\times10^{-8}$~см & см & \eqref{eq:thomas-fermi} \\
\end{longtblr}

\par\Needspace{7\baselineskip}\medskip\noindent\textbf{Колективність і зв'язаність}\par\nobreak\smallskip
\noindent\begin{longtblr}{colspec={X[0.9,c,m]X[3.4,l,m]X[1.5,c,m]X[1.5,c,m]}}
Символ & Значення & Одиниця (СГС) & Де \\
$N_D$ & число електронів у~дебаївській сфері $\frac{4\pi}{3}n_e\lambda_{De}^3$ & --- & \eqref{eq:ND} \\
$a$ & середня відстань між частинками $n^{-1/3}$ & см & §\,\ref{sec:debye-sphere}\,(с.\,\pageref{sec:debye-sphere}) \\
$U_{\mathrm C}$, $U_T$ & енергія кулонівської взаємодії $e^2/a$ та теплова енергія $T$ & еВ & §\,\ref{sec:coupling}\,(с.\,\pageref{sec:coupling}) \\
$\Gamma$ & параметр зв'язку $e^2/(aT)$ & --- & \eqref{eq:Gamma} \\
$p$ & тиск ідеального газу $\sum_sn_sT_s$ & дин$/$см$^{2}$ & §\,\ref{sec:coupling}\,(с.\,\pageref{sec:coupling}) \\
\end{longtblr}

\par\Needspace{7\baselineskip}\medskip\noindent\textbf{Шар біля стінки та зонд Ленгмюра}\par\nobreak\smallskip
\noindent\begin{longtblr}{colspec={X[0.9,c,m]X[3.4,l,m]X[1.5,c,m]X[1.5,c,m]}}
Символ & Значення & Одиниця (СГС) & Де \\
$\sigma$ & поверхнева густина заряду площини & СГСЕ$/$см$^{2}$ & §\,\ref{sec:debye-layer}\,(с.\,\pageref{sec:debye-layer}) \\
$x$, $y$ & відстань від площини; $y=e\varphi/T$ & см; --- & §\,\ref{sec:debye-layer}\,(с.\,\pageref{sec:debye-layer}) \\
$c_s$ & іонно-звукова швидкість $\sqrt{T_e/m_i}$ при $T_i\ll T_e$ (загальна~--- розд.~\ref{sec:mhd-waves}) & см/с & \eqref{eq:bohm-criterion} \\
$u$, $u_0$ & швидкість іонів у~шарі; їхня швидкість на вході в~шар (критерій Бома $u_0\ge c_s$) & см/с & \eqref{eq:bohm-criterion} \\
$n_0$, $n_\mathrm{sh}$ & концентрація в~незбуреній плазмі; на межі шару ($n_\mathrm{sh}\approx0{,}6\,n_0$) & см$^{-3}$ & §\,\ref{sec:bohm-sheath}\,(с.\,\pageref{sec:bohm-sheath}) \\
$\chi$ & безрозмірний потенціал $-e\varphi/T_e$ & --- & §\,\ref{sec:bohm-sheath}\,(с.\,\pageref{sec:bohm-sheath}) \\
$\Gamma_i$, $\Gamma_e$ & потік іонів, електронів на стінку & см$^{-2}\cdot$с$^{-1}$ & §\,\ref{sec:bohm-sheath}\,(с.\,\pageref{sec:bohm-sheath}) \\
$\bar v_e$, $v_{Te}$, $v_{Ti}$ & середня швидкість електрона $\sqrt{8T_e/\pi m_e}$; теплові швидкості електронів і~іонів & см/с & §\,\ref{sec:bohm-sheath}\,(с.\,\pageref{sec:bohm-sheath}) \\
$\varphi_f$ & плаваючий потенціал відносно плазми, $\varphi_f=V_f-V_p<0$ & СГС (потенціал) & \eqref{eq:floating-potential} \\
$V$, $V_f$, $V_p$ & потенціал зонда відносно опорного електрода; плаваючий потенціал; потенціал плазми & В & §\,\ref{sec:bohm-sheath}\,(с.\,\pageref{sec:bohm-sheath}) \\
$I$ & струм на зонд (не плутати з~енергією іонізації $I$) & мА & §\,\ref{sec:bohm-sheath}\,(с.\,\pageref{sec:bohm-sheath}) \\
$I_{is}$, $I_{es}$, $I_e$ & струми іонного та електронного насичення; електронний струм & мА & §\,\ref{sec:bohm-sheath}\,(с.\,\pageref{sec:bohm-sheath}) \\
$S$, $r_p$ & площа зонда; його характерний розмір & см$^{2}$; см & §\,\ref{sec:bohm-sheath}\,(с.\,\pageref{sec:bohm-sheath}) \\
$V_1$, $V_2$, $I_{e1}$, $I_{e2}$ & потенціали й~електронні струми двох електродів подвійного зонда & В; мА & \eqref{eq:double-probe} \\
$\lambda_{\mathrm{fp}}$ & довжина вільного пробігу (відстань між зіткненнями) & см & §\,\ref{sec:bohm-sheath}\,(с.\,\pageref{sec:bohm-sheath}) \\
$\vect B$, $B$ & індукція магнітного поля & Гс & \eqref{eq:larmor} \\
$v_\perp$ & складова швидкості поперек $\vect B$ & см/с & \eqref{eq:larmor} \\
$r_{Le}$ & ларморівський радіус електрона $m_ecv_\perp/eB$ (розд.~\ref{sec:particle-motion}) & см & \eqref{eq:larmor} \\
\end{longtblr}

\par\Needspace{7\baselineskip}\medskip\noindent\textbf{Хвилі в плазмі (розд.~\ref{sec:plasma-waves})}\par\nobreak\smallskip
\noindent\begin{longtblr}{colspec={X[0.9,c,m]X[3.4,l,m]X[1.5,c,m]X[1.5,c,m]}}
Символ & Значення & Одиниця (СГС) & Де \\
$\omega_{pe}$, $\omega_p$ & електронна плазмова частота $\sqrt{4\pi n_ee^2/m_e}$; далі $\omega_p\equiv\omega_{pe}$ & с$^{-1}$ & \eqref{eq:plasma-frequency} \\
$\varepsilon$ & діелектрична проникність $1-\omega_p^2/\omega^2$ & --- & \eqref{eq:plasma-permittivity} \\
$N$ & показник заломлення $\sqrt\varepsilon$ (літеру $n$ зарезервовано для концентрації) & --- & §\,\ref{sec:plasma-permittivity}\,(с.\,\pageref{sec:plasma-permittivity}) \\
$v_\phi$, $v_g$ & фазова $\omega/k$ та групова $d\omega/dk$ швидкості & см/с & §\,\ref{sec:plasma-permittivity}\,(с.\,\pageref{sec:plasma-permittivity}) \\
$n_{\mathrm{кр}}$ & критична концентрація $m_e\omega^2/4\pi e^2$, при якій $\omega=\omega_p$ & см$^{-3}$ & \eqref{eq:critical-density} \\
$\nu$ & ефективна частота зіткнень електрона ($\nu_{ei}$) & с$^{-1}$ & \eqref{eq:eps-collisions-waves} \\
$\delta$ & товщина скін-шару, глибина проникнення $1/\kappa$ & см & §\,\ref{sec:plasma-permittivity}\,(с.\,\pageref{sec:plasma-permittivity}) \\
\end{longtblr}

\par\Needspace{7\baselineskip}\medskip\noindent\textbf{Хвилі в замагніченій плазмі (розд.~\ref{sec:mhd-waves})}\par\nobreak\smallskip
\noindent\begin{longtblr}{colspec={X[0.9,c,m]X[3.4,l,m]X[1.5,c,m]X[1.5,c,m]}}
Символ & Значення & Одиниця (СГС) & Де \\
$\vect B_0$, $\theta$ & рівноважне поле (вздовж осі $z$); кут між $\vect k$ і~$\vect B_0$ & Гс; рад & \eqref{eq:magnetosonic-biquad} \\
$\rho_m$, $\rho_0$, $\delta\rho$ & масова густина, її рівноважне значення й~збурення (не плутати з~густиною заряду $\rho$) & г$/$см$^{3}$ & \eqref{eq:wave-linear-continuity} \\
$p_e$, $p_i$, $\beta$ & тиск електронів, іонів; $\beta=8\pi p/B^2$ & дин$/$см$^{2}$; --- & розд.~\ref{sec:mhd} \\
$\gamma$, $\gamma_e$, $\gamma_i$ & показник адіабати: спільний; електронів ($\gamma_e=1$, ізотермічні); іонів ($\gamma_i=3$, одновимірне стискання) & --- & §\,\ref{sec:sound-wave}\,(с.\,\pageref{sec:sound-wave}) \\
$c_s$ & швидкість звуку $\sqrt{\gamma(p_e+p_i)/\rho_m}$; для іонного звуку $\sqrt{(ZT_e+3T_i)/m_i}$ & см/с & \eqref{eq:wave-sound-equation}, \eqref{eq:ion-acoustic} \\
$v_A$, $\vect v_A$ & альвенівська швидкість $B_0/\sqrt{4\pi\rho_0}$ (вектор вздовж $\vect B_0$) & см/с & \eqref{eq:wave-alfven-speed} \\
$v_\pm$ & фазові швидкості швидкої ($+$) і~повільної ($-$) магнітозвукових мод & см/с & \eqref{eq:wave-magnetosonic-general} \\
$\nu_m$ & магнітна в'язкість $c^2/4\pi\sigma$ (коефіцієнт дифузії поля) & см$^{2}/$с & §\,\ref{sec:linear-mhd}\,(с.\,\pageref{sec:linear-mhd}) \\
$\omega_{ce}$, $\omega_{ci}$, $\Omega_s$ & циклотронні частоти електрона $eB_0/m_ec$ та~іона $ZeB_0/m_ic$; $\Omega_s=q_sB_0/m_sc$ зі~знаком & с$^{-1}$ & §\,\ref{sec:cold-magnetized}\,(с.\,\pageref{sec:cold-magnetized}) \\
$\omega_{ps}$ & плазмова частота сорту $s$ ($\omega_{pe}$, $\omega_{pi}$) & с$^{-1}$ & \eqref{eq:stix-RLP} \\
$\hat\varepsilon$; $S$, $D$, $P$, $R$, $L$ & тензор діелектричної проникності та~його елементи (позначення Стікса) & --- & \eqref{eq:stix-tensor}, \eqref{eq:stix-RLP} \\
$\omega_R$, $\omega_L$ & частоти відсікання $R$- і~$L$-хвиль вздовж поля & с$^{-1}$ & §\,\ref{sec:cold-magnetized}\,(с.\,\pageref{sec:cold-magnetized}) \\
$\omega_{UH}$, $\omega_{LH}$ & верхній і~нижній гібридні резонанси & с$^{-1}$ & \eqref{eq:upper-hybrid} \\
$\psi$ & кут обертання Фарадея & рад & \eqref{eq:faraday-rotation} \\
\end{longtblr}
}

\paragraph{Символи з кількома значеннями.} У~розд.~\ref{sec:plasma-basics} і~\ref{sec:mhd-waves} деякі літери означають різні величини;
значення визначається контекстом та індексами.
\begin{itemize}
	\item $I$~--- енергія іонізації (\S\,\ref{sec:saha}) і~струм на зонд (\S\,\ref{sec:bohm-sheath});
	\item $\Gamma$~--- параметр зв'язку~\eqref{eq:Gamma} ; з~індексом виду ($\Gamma_i$, $\Gamma_e$)~--- потік частинок (на стінку, при дифузії);
	\item $n_s$~--- концентрація сорту $s$ (§\,\ref{sec:debye-screening}\,(с.\,\pageref{sec:debye-screening})); концентрацію на межі шару позначено $n_\mathrm{sh}$;
	\item $a$~--- середня відстань між частинками (розмір зонда позначено $r_p$);
	\item $\rho$~--- густина заряду; ларморівський радіус позначено $r_{L}$ ($r_{Le}$ для електрона), а~не $\rho$;
	\item $u$~--- швидкість іонів у~шарі; допоміжна змінна $u=r\varphi$ у~§\,\ref{sec:debye-screening}\,(с.\,\pageref{sec:debye-screening}) вживається лише в~одному виведенні;
	\item $\lambda$ без індексу~--- довжина екранування в~загальному випадку~\eqref{eq:screening-general}, $\lambda_D$~--- дебаївська;
	\item $E$~--- напруженість електричного поля, $E_F$~--- енергія Фермі;
	\item $N$~--- показник заломлення (розд.~\ref{sec:plasma-waves}); у~розд.~\ref{sec:plasma-basics} $N_D$~--- число частинок у~дебаївській сфері, а~$n$~--- завжди концентрація;
	\item $\delta$~--- товщина скін-шару (розд.~\ref{sec:plasma-waves}) і~дельта-функція $\delta(\vect r)$;
	\item $S$~--- площа зонда (розд.~\ref{sec:plasma-basics}) і~елемент тензора Стікса (розд.~\ref{sec:mhd-waves});
	\item $\lambda$~--- довжина екранування (розд.~\ref{sec:plasma-basics}) і~довжина хвилі в~обертанні Фарадея (розд.~\ref{sec:mhd-waves});
	\item $c_s$~--- $\sqrt{T_e/m_i}$ у~розд.~\ref{sec:plasma-basics} (межа $T_i\ll T_e$) і~загальна швидкість звуку в~розд.~\ref{sec:mhd-waves};
	\item $\rho_m$~--- масова густина (розд.~\ref{sec:mhd-waves}), на~відміну від густини заряду $\rho$.
\end{itemize}

